\documentclass[11pt, a4paper]{article}

\usepackage[margin=2.5cm]{geometry}
\usepackage{amsmath, amssymb, amsthm}
\usepackage{booktabs}
\usepackage{array}
\usepackage{microtype}
\usepackage{hyperref}
\usepackage{xcolor}
\usepackage{parskip}
\usepackage{titlesec}

\hypersetup{colorlinks=true, linkcolor=black, citecolor=black, urlcolor=blue}

\titleformat{\section}{\large\bfseries}{}{0em}{}[\titlerule]
\titleformat{\subsection}{\normalsize\bfseries}{}{0em}{}

\setlength{\parskip}{6pt}

\title{\textbf{Differential Inverse Kinematics}\\
       \large Mathematical Foundations of the xArm7 Control Pipeline}
\author{}
\date{}

% -----------------------------------------------------------------------
\begin{document}
\maketitle
\thispagestyle{empty}

% -----------------------------------------------------------------------
\section{Problem Statement}

A robot manipulator with $n$ joints has a \emph{configuration} $\mathbf{q} \in \mathbb{R}^n$ (the vector of joint angles) and an \emph{end-effector (EE) pose} $\mathbf{x} \in \mathbb{R}^m$ in task space. For the xArm7, $n = 7$ and $\mathbf{x}$ is the 3-D Cartesian position of the EE site, so $m = 3$.

Forward kinematics (FK) maps configuration to pose:
\[
  \mathbf{x} = f(\mathbf{q})
\]

\emph{Inverse kinematics} (IK) is the inverse problem: given a desired EE motion $\Delta \mathbf{x}$, find the joint displacement $\Delta \mathbf{q}$ that produces it. Because $f$ is nonlinear and the system is redundant ($n > m$), a closed-form solution is generally intractable. Instead, the \emph{differential} (velocity-level) formulation is used.

% -----------------------------------------------------------------------
\section{The Jacobian Matrix}

Differentiating FK with respect to time gives the \emph{Jacobian} $\mathbf{J}(\mathbf{q})$, a $m \times n$ matrix that linearly relates joint velocities to EE velocity:
\[
  \dot{\mathbf{x}} = \mathbf{J}(\mathbf{q})\,\dot{\mathbf{q}}, \qquad
  J_{ij}(\mathbf{q}) = \frac{\partial x_i}{\partial q_j}
\]

The Jacobian is configuration-dependent and must be recomputed at each control step. In the simulation it is evaluated using MuJoCo's \texttt{mj\_jacSite}, which returns the $3 \times 7$ \emph{translational} Jacobian at the EE site:
\[
  \mathbf{J} = \frac{\partial \mathbf{p}_{\text{EE}}}{\partial \mathbf{q}} \in \mathbb{R}^{3 \times 7}
\]

where $\mathbf{p}_{\text{EE}} \in \mathbb{R}^3$ is the Cartesian position of the EE site. Orientation is not included because the teleop interface commands only translational deltas.

% -----------------------------------------------------------------------
\section{The Naive Pseudo-Inverse}

Given a desired EE displacement $\Delta \mathbf{x}$, the IK problem at the differential level is:
\[
  \Delta \mathbf{x} = \mathbf{J}\,\Delta \mathbf{q}
\]

This is an underdetermined linear system ($m < n$). The minimum-norm least-squares solution is given by the Moore-Penrose pseudo-inverse:
\[
  \Delta \mathbf{q} = \mathbf{J}^{+}\,\Delta \mathbf{x}, \qquad
  \mathbf{J}^{+} = \mathbf{J}^T\!\left(\mathbf{J}\mathbf{J}^T\right)^{-1}
\]

This selects the solution $\Delta \mathbf{q}$ with smallest Euclidean norm, which tends to distribute motion evenly across joints and avoid unnecessary movement in the null space.

% -----------------------------------------------------------------------
\section{Singularities and the Damped Least Squares Method}

\subsection{The Singularity Problem}

A kinematic \emph{singularity} occurs when $\mathbf{J}$ loses full row rank — equivalently, when $\det(\mathbf{J}\mathbf{J}^T) \approx 0$. Physically, this corresponds to configurations where the arm is fully extended, folded back on itself, or aligned along a degenerate axis. Near such configurations:
\[
  \left(\mathbf{J}\mathbf{J}^T\right)^{-1} \longrightarrow \infty
\]

The naive pseudo-inverse amplifies small EE errors into arbitrarily large joint velocities, causing the simulated arm to \emph{jitter} or \emph{explode} at workspace boundaries.

\subsection{Levenberg-Marquardt Regularisation}

The \emph{damped least-squares} (DLS) method, also known as Levenberg-Marquardt regularisation, adds a damping term $\lambda^2\mathbf{I}$ to the matrix that must be inverted:
\[
  \boxed{
    \mathbf{J}_{\text{damp}} = \mathbf{J}^T\!\left(\mathbf{J}\mathbf{J}^T + \lambda^2\mathbf{I}\right)^{-1}
  }
\]

The regularised joint displacement is then:
\[
  \Delta \mathbf{q} = \mathbf{J}_{\text{damp}}\,\Delta \mathbf{x}
\]

\subsection{Effect of the Damping Factor}

The matrix $(\mathbf{J}\mathbf{J}^T + \lambda^2\mathbf{I})$ has eigenvalues $\sigma_i^2 + \lambda^2 > 0$ for all singular values $\sigma_i \geq 0$, so it is always invertible regardless of the arm configuration. This can be understood via the SVD of $\mathbf{J} = \mathbf{U}\mathbf{\Sigma}\mathbf{V}^T$:

\[
  \mathbf{J}_{\text{damp}} = \mathbf{V}\,\tilde{\mathbf{\Sigma}}\,\mathbf{U}^T, \qquad
  \tilde{\sigma}_i = \frac{\sigma_i}{\sigma_i^2 + \lambda^2}
\]

When $\sigma_i \gg \lambda$ (away from singularity), $\tilde{\sigma}_i \approx \sigma_i^{-1}$ and the solution approaches the pseudo-inverse. When $\sigma_i \ll \lambda$ (near singularity), $\tilde{\sigma}_i \approx \sigma_i / \lambda^2$ and the contribution of that mode is suppressed rather than amplified. The trade-off is accuracy: DLS introduces a small EE tracking error proportional to $\lambda^2 / \sigma_i$ in the affected directions.

In the implementation the damping factor is fixed at $\lambda = 0.05$, providing robustness without noticeably degrading tracking under normal operation.

% -----------------------------------------------------------------------
\section{Motion Smoothing}

Commands from the keyboard arrive as discrete impulses (a fixed $\Delta \mathbf{x}_{\text{step}}$ per keypress). Applied directly to a 1000 Hz physics engine, these would create high-frequency force spikes. A first-order exponential filter smooths the motion profile.

\subsection{Exponential Decay Filter}

Commands accumulate in a \emph{move queue} $\mathbf{m} \in \mathbb{R}^3$. At each physics timestep $\Delta t$ the queue is drained by an exponentially decaying fraction:
\[
  \alpha = 1 - e^{-k\,\Delta t}
\]
\[
  \Delta \mathbf{x}_{\text{applied}} = \alpha\,\mathbf{m}, \qquad
  \mathbf{m} \leftarrow (1 - \alpha)\,\mathbf{m} = \mathbf{m} - \Delta \mathbf{x}_{\text{applied}}
\]

With $k = 10.0$ and $\Delta t = 0.001\,\text{s}$, the filter time constant is $\tau = 1/k = 0.1\,\text{s}$, giving a smooth exponential approach to each target that damps out oscillations in the MuJoCo solver.

The peak EE speed from a single queued command of magnitude $\|\mathbf{m}\|$ is:
\[
  v_{\text{peak}} = \frac{\|\mathbf{m}\|\,\alpha}{\Delta t} \approx \|\mathbf{m}\| \cdot k
\]

\subsection{Queue Cap}

To prevent prolonged motion after key-release (backlog) the queue norm is capped before draining:
\[
  \mathbf{m} \leftarrow \mathbf{m}\cdot\min\!\left(1,\,\frac{v_{\max}}{k\,\|\mathbf{m}\|}\right)
\]

where $v_{\max} = 0.5\,\text{m/s}$. This bounds the instantaneous EE speed at all times.

% -----------------------------------------------------------------------
\section{Speed Limiting}

\subsection{Cartesian Speed Limit}

The queue cap described above ensures:
\[
  \left\|\dot{\mathbf{p}}_{\text{EE}}\right\|_2 \leq v_{\max} = 0.5\,\text{m/s}
\]

\subsection{Joint Speed Limit}

After computing $\Delta \mathbf{q}$, an additional per-joint speed check is applied. The joint velocity at step $k$ is $\dot{q}_i = \Delta q_i / \Delta t$. The entire vector is scaled uniformly to ensure no joint exceeds $\omega_{\max} = \pi\,\text{rad/s}$ (the xArm7 rated maximum):
\[
  s = \frac{\omega_{\max}\,\Delta t}{\max_i |\Delta q_i| + \epsilon}, \qquad
  \Delta \mathbf{q} \leftarrow \min(s,\,1)\cdot\Delta \mathbf{q}
\]

Scaling the whole vector uniformly preserves the direction of EE motion; only the speed is reduced. The small $\epsilon = 10^{-9}$ avoids division by zero.

% -----------------------------------------------------------------------
\section{Full Control Loop}

The complete per-timestep procedure is summarised below.

\begin{center}
\renewcommand{\arraystretch}{1.4}
\begin{tabular}{>{\bfseries}r l}
\toprule
Step & Operation \\
\midrule
1 & Receive teleop delta $\Delta\mathbf{x}_{\text{cmd}}$; accumulate into $\mathbf{m}$. \\
2 & Cap: $\mathbf{m} \leftarrow \mathbf{m}\cdot\min\!\left(1,\; v_{\max}/\left(k\,\|\mathbf{m}\|\right)\right)$. \\
3 & Filter: $\Delta\mathbf{x} = \alpha\,\mathbf{m}$;\; $\mathbf{m} \leftarrow \mathbf{m} - \Delta\mathbf{x}$. \\
4 & Compute Jacobian $\mathbf{J}$ at current $\mathbf{q}$ via \texttt{mj\_jacSite}. \\
5 & Solve: $\Delta\mathbf{q} = \mathbf{J}^T\!\left(\mathbf{J}\mathbf{J}^T + \lambda^2\mathbf{I}\right)^{-1}\Delta\mathbf{x}$. \\
6 & Clamp: $\Delta\mathbf{q} \leftarrow \min\!\left(1,\; \omega_{\max}\Delta t / \max_i|\Delta q_i|\right)\cdot\Delta\mathbf{q}$. \\
7 & Update: $\mathbf{q}_{\text{ctrl}} \leftarrow \text{clip}\!\left(\mathbf{q}_{\text{ctrl}} + \Delta\mathbf{q},\; \mathbf{q}_{\min}, \mathbf{q}_{\max}\right)$. \\
8 & Advance simulation: \texttt{mj\_step}$(\Delta t = 1\,\text{ms})$. \\
\bottomrule
\end{tabular}
\end{center}

% -----------------------------------------------------------------------
\section{Parameter Summary}

\begin{center}
\renewcommand{\arraystretch}{1.3}
\begin{tabular}{lll}
\toprule
Parameter & Value & Role \\
\midrule
$n$ & 7 & Number of arm joints (xArm7) \\
$m$ & 3 & Task-space dimension (Cartesian position) \\
$\lambda$ & $0.05$ & LM damping factor \\
$k$ & $10.0\,\text{s}^{-1}$ & Exponential filter gain \\
$\Delta t$ & $0.001\,\text{s}$ & Physics timestep (1000 Hz) \\
$v_{\max}$ & $0.5\,\text{m/s}$ & Maximum EE Cartesian speed \\
$\omega_{\max}$ & $\pi\,\text{rad/s}$ & Maximum joint speed (180\,°/s) \\
\bottomrule
\end{tabular}
\end{center}

\end{document}
